\section{Introduction}
\label{sec:introduction}

WireGuard decrypts a tunnel's packets in parallel to raise throughput, yet must
hand them to the network stack in the order they arrived. On the receive path,
each data packet is recorded in a per-peer queue in arrival order and, at the
same time, published to a device-wide ring that per-CPU workers drain and
decrypt~\cite{donenfeld2017,wireguard-linux}. Parallel decryption is what lets a
single tunnel use more than one core. Ordered delivery is a protocol
requirement: the packet due next must be committed before any packet behind it.
These two requirements meet at one difficult point. The packet due for delivery
next may be one of the last to finish decrypting.

Consider two packets from one tunnel. \texttt{P1} arrives first and becomes the
ordered head; \texttt{P2} arrives behind it. Both are dispatched to the shared
decrypt ring, and a worker completes \texttt{P2} before the worker handling
\texttt{P1}. \texttt{P2} is now decrypted, but it cannot be delivered: the
receive poll advances the per-peer queue only from its head, and the head,
\texttt{P1}, is still encrypted. The poll reaches \texttt{P1}, stops, and
delivers nothing, even though ready data waits behind it and pending decrypt
work may still sit in the shared ring. We call this an execution-order
inversion. The consumer that would deliver is runnable, but the work it needs is
owned by a different scheduling path. Ordered delivery itself is correct here;
the inefficiency is that the blocked context waits while useful work remains
available.

This inversion matters most for a single tunnel. On a multicore receiver, the
NIC hashes each packet's flow to a receive queue, and distinct flows land on
distinct cores~\cite{linux-napi,salim2001beyond}. Flow hashing therefore spreads
\emph{many} tunnels across cores, but it does not divide \emph{one}: a WireGuard
tunnel is a single outer UDP flow with fixed addresses and ports, so its packets
hash to one receive queue and its ordered delivery runs in one RX/NAPI context.
In our single-tunnel measurements, that context's softirq processing occupies
close to one full core and becomes the throughput ceiling. Site-to-site and
single-heavy-tunnel deployments stay bound to one receive core no matter how
many cores the machine has, and the inversion is felt exactly there.

Because a completed non-head packet wakes the receive poll, an inversion
produces polls that run before the head is ready and deliver nothing.
Suppressing those premature wakes is a natural first step, and one that connects
to prior work on the same receive path, which improved it by changing where
deferred receive processing executes~\cite{mounah2025}. We implemented two
composable wake-side changes that remove premature wakes directly. They reduce
the wasted-poll fraction substantially, from about 27\% to 13--15\% across peer
counts, and their effect strengthens as decryption is made slower, which
confirms that they act on the inversion. Direct accounting, however, shows the
reclaimed work is small: the entire population of wasted polls occupies only
about two hundredths of a core on the tested hardware. The wake-side changes
remove many events, but each event is cheap, so no throughput or aggregate-CPU
benefit follows on this platform. The changes are real and they work as
intended; they are simply not where the cost lives.

The quantity that matters is not how often the poll fires early, but how long
delivery stays blocked behind an encrypted head. Classified instrumentation
shows those intervals last tens of microseconds: the head is not slow to
decrypt, it is slow to be \emph{scheduled} for decryption, while the blocked
context sits idle. \texttt{wg\_steal} puts that idle context to work. When
ordered delivery stalls behind an incomplete head, the blocked NAPI context
performs a bounded amount of pending decrypt work itself before yielding. It
consumes jobs from the device-wide decrypt ring, decrypts them on its own core,
and publishes each through WireGuard's normal completion path, then rechecks the
head and resumes ordered delivery. Normal workers keep running in parallel, the
borrowed work is bounded, and packet order is never bypassed.

% EVIDENCE REVIEW: the headline figures restated in this paragraph
% (+1.96% throughput and -3.66% total CPU for the saturated confirmation;
% -1.17% CPU non-detection for the matched-load confirmation) are taken
% verbatim from docs/paper/CLAIM_EVIDENCE_MATRIX.md (C11, C12, C17) and
% paper/tables/confirmation_results.tex. Restated for the reader, not recomputed.
We evaluate a bounded budget of four jobs per stealing pass in two paired
confirmations on 10-GbE CloudLab hardware. In the uncapped saturated
single-tunnel regime, where the receive core is the binding resource, stealing
raised throughput by 1.96\% and reduced total busy CPU by 3.66\%; both were
declared co-primary endpoints, and both moved in the favorable direction. At an
exactly matched delivered load near 3.8~Gb/s, no favorable CPU effect was
detected: the estimate is $-1.17\%$, with a confidence interval that spans both
a saving and a small cost. Mechanism counters show that the classified
blocked-head population was still nearly eliminated at that load, so the
mechanism remained active even where its aggregate-CPU effect was not
detectable. This pattern is consistent with stealing becoming measurable as the
pinned receive core approaches its binding limit. It does not establish a
saturation threshold, and we make no user-visible latency claim.

This paper makes the following contributions.
\begin{itemize}
  \item A source-level characterization of execution-order inversion and
    ordered-head blocking in WireGuard's parallel receive path
    (Section~\ref{sec:background}, Section~\ref{sec:design-model}).
  \item Instrumentation that separates premature polling from longer-lived
    encrypted-head waiting, and that quantifies the cost of each
    (Section~\ref{sec:eval-poll-cost}, Section~\ref{sec:eval-blocking}).
  \item \texttt{wg\_steal}, a bounded, ordering-preserving mechanism that lets a
    blocked receive context perform pending decrypt work from the device-wide
    ring (Section~\ref{sec:design-steal}).
  \item A paired CloudLab evaluation that detects a throughput and total-CPU
    benefit in the uncapped saturated regime and reports a valid matched-load
    CPU non-detection near 3.8~Gb/s (Section~\ref{sec:eval-gatea},
    Section~\ref{sec:eval-gateb}).
  \item An explicit account of the mechanism's limits, including fairness and
    softirq-duration trade-offs and the absence of a valid user-visible latency
    result (Section~\ref{sec:design-limitations},
    Section~\ref{sec:eval-limitations}).
\end{itemize}
